\section{引言}

循证医学问答要求系统在生成或选择答案之前检索到相关的生物医学证据。检索增强生成（RAG）为结合检索证据与语言模型提供了通用框架 \citep{lewis2020rag}，近期医学 RAG 基准研究也表明，检索组件选择会显著影响医学问答效果 \citep{xiong2024benchmarking}。然而，常规词法检索或稠密检索通常只建模问题与段落之间的成对相似度，而生物医学证据常常由问题、段落、MeSH 概念、生物医学实体、文档和关系来源之间的高阶关系共同组织。

本文研究的是一个实践性、领域特定的重排序方法，而不是新的通用 RAG 范式。我们提出 KCH-MedRank：一种知识约束的局部超图学习排序框架，用于生物医学证据检索与重排序。该方法使用 BM25、稠密检索和倒数排名融合保持候选生成的可复现性 \citep{robertson2009bm25,cormack2009rrf}，随后围绕每个问题及其候选证据段落构建局部跨粒度超图。

本文贡献定位保持克制并以证据为中心：提出一个轻量重排序模块，结合生物医学语义信号、MeSH 层级感知特征、局部超图扩散与中心性特征，以及查询分组 LambdaMART 排序，以提升证据检索效果。该系统不声称替代临床判断，不提供临床诊断，也不把答案生成作为主要贡献。

当前 BioASQ 证据检索实验主线表明，增强候选生成提高了留出测试集 top-100 召回上限，KCH-MedRank 相对于原始 Hybrid RRF 和真实 MedCPT Cross-Encoder 基线提升了 top-10 证据检索效果。失败分析也显示存在丢失的金标准证据，因此讨论部分将该方法表述为具有明确净收益但仍存在失败模式的证据检索方法。

